\section{Discussion} \label{sec:5}

\subsection{On BPSO and BBOA's Early-Generation Fitness Advantage} \label{sec:5.1}
Figures \ref{fig:4.1.1} and \ref{fig:4.2.1} show BPSO and BBOA reaching a markedly better mean best-fitness than every other method within the very first recorded evaluations, before plateauing early and eventually being overtaken by BGWO, HybridGWO, and the proposed CSO pipeline. This is not the result of an informed or non-random initialization scheme: every algorithm in this study, including BPSO and BBOA, draws its initial population from the same uniform distribution with no warm-start, opposition-based, or chaotic-map seeding. The advantage instead comes from the hybrid S-shaped/V-shaped transfer function that Hashemi et al. \cite{Hashemi:2026} specify for BPSO and BBOA (Section \ref{sec:3.7}): for every individual, at every generation including the very first, the same continuous position is binarized twice, once under the S-shaped rule of Eq. (\ref{eq:s-shaped-baseline}) and once under the V-shaped rule of Eq. (\ref{eq:v-shaped}), both candidates are evaluated, and whichever achieves the lower fitness is kept. Because a best-of-two selection cannot score worse, in expectation, than a single draw, BPSO and BBOA's reported fitness at any generation is a best-of-two order statistic rather than a single sample from the same distribution the other baselines and the proposed method draw from, which explains why their initial population already looks stronger despite an identical, uninformed random start.

This mechanism also means BPSO and BBOA spend two fitness evaluations per individual per generation wherever the other five methods spend one, so under the shared evaluation budget $FE_{max}$ (Section \ref{sec:3.7}) they complete roughly half as many generations. Their early lead is therefore a one-time gain from the transfer function's built-in selection pressure rather than evidence of a more effective search or exploitation mechanism; it does not compound across generations the way the winner/loser competitive update of CSO (Eq. (\ref{eq:3.6.1})) or the GWO family's convergence-controlled encircling does, which is consistent with BPSO and BBOA's curves flattening earliest in Figures \ref{fig:4.1.1} and \ref{fig:4.2.1} while every other method continues to improve.

\subsection{On the Comparison Between CSO and the Baselines} \label{sec:5.2}
CSO's advantage over the five baselines is dataset-dependent rather than uniform. On Naranjo, it attains the highest mean balanced accuracy of all six methods (0.847, versus 0.827--0.843 for the baselines; Figure \ref{fig:4.2.2}), whereas on Oxford it ranks second (0.810), behind HybridGWO's 0.815 (Figure \ref{fig:4.1.2}). A paired Wilcoxon signed-rank test on balanced accuracy, matched by seed across algorithms ($n=20$ per comparison), confirms this asymmetry. On Oxford, CSO significantly outperforms BGWO ($p=0.048$), QMFO ($p=0.007$), and BPSO ($p=0.036$), and is on the margin against BBOA ($p=0.050$); only its comparison with QMFO survives a Bonferroni-corrected threshold of $\alpha=0.01$ for the five simultaneous comparisons. Against HybridGWO, CSO is statistically indistinguishable ($p=0.51$, a near-even 10/20 pairwise split) -- the only baseline not clearly outperformed on either dataset. On Naranjo, CSO significantly outperforms all five baselines (BGWO $p=0.028$, HybridGWO $p=0.079$, QMFO $p<0.001$, BPSO $p<0.001$, BBOA $p<0.001$), of which the comparisons against QMFO, BPSO, and BBOA remain significant under the same Bonferroni-corrected threshold.

This pattern is consistent with the higher dimensionality of Naranjo's encoding: at $D=45$ acoustic features, versus $D=22$ for Oxford, the joint search space of Section \ref{sec:3.3} is roughly twice as large, giving CSO's per-particle searched transfer function -- letting each candidate solution pick whichever of five binarization rules (Section \ref{sec:3.6.2}) suits its own region of the space, rather than every particle sharing one fixed rule -- correspondingly more room to specialize. On the smaller Oxford space, this adaptive advantage is thinner, and HybridGWO's fixed WOA-to-GWO cascade schedule (Table \ref{tab:hyperparameters}) is competitive without it.

CSO's efficiency profile (Figures \ref{fig:4.1.4} and \ref{fig:4.2.4}) also differs by dataset. On Oxford, it selects the smallest mean feature ratio of all six methods (0.225, versus 0.227--0.300 for the baselines) while simultaneously improving accuracy, that is, no accuracy-parsimony trade-off is paid. On Naranjo, it selects more features on average (0.261) than BBOA (0.177) or BPSO (0.234), trading some parsimony for its accuracy lead. Because the fitness function (Eq. (\ref{eq:3.5.1})) weights balanced accuracy nine times as heavily as feature ratio ($\omega=0.9$), this is the expected behavior rather than an inconsistency: whenever a larger feature subset yields a material accuracy gain, the search is designed to keep it. QMFO selects the largest feature subsets on both datasets (0.300 on Oxford, 0.377 on Naranjo) without a corresponding accuracy advantage, suggesting its search dynamics are comparatively less effective than the other methods at pruning uninformative features under this joint encoding.

\subsection{Clinical Implications} \label{sec:5.3}
The balanced accuracies attained by every method in this comparison, roughly 0.79--0.85 across both datasets (Tables \ref{tab:results-oxford} and \ref{tab:results-naranjo}), are consistent with what voice-based machine-learning screening for PD has generally reported in the wider literature \cite{SedighMalekroodi:2025}: useful enough to support a low-cost, non-invasive triage or monitoring tool, but well short of a diagnostic-grade instrument that could substitute for the MDS Clinical Diagnostic Criteria discussed in Section \ref{sec:2.1}. The intended clinical role of a pipeline like this one is therefore as a first-pass flag, prompting referral for clinical evaluation, or as a between-visit monitoring signal in telemedicine and home-recording settings where repeated in-person motor examination is impractical, rather than as a standalone diagnostic decision-maker.

The subject-grouped cross-validation protocol of Section \ref{sec:3.2} is not a minor implementation detail from this clinical perspective: because Oxford's subjects contribute an average of six recordings each and Naranjo's contribute exactly three, a record-level split would let recordings from the same speaker appear in both the training and validation folds, letting the classifier partly learn speaker identity rather than PD status, which would inflate every metric in Tables \ref{tab:results-oxford} and \ref{tab:results-naranjo} relative to a pipeline's real accuracy on a genuinely new patient. Every result reported in this study instead estimates accuracy on speakers unseen during that fold's training, the quantity that actually matters for deployment on a new patient rather than a new recording from an already-seen one.

Finally, the dataset-dependent ranking discussed in Section \ref{sec:5.2}, where HybridGWO is competitive on Oxford but not on Naranjo, and CSO's advantage strengthens with the higher-dimensional Naranjo encoding, cautions against treating any single method, including the one proposed here, as universally best for voice-based PD screening. A practitioner adapting this pipeline to a new population or recording protocol should budget for re-running this kind of multi-method comparison on their own data rather than assuming the ranking observed on Oxford and Naranjo transfers unchanged.

\subsection{Limitations and Future Work} \label{sec:5.4}
As noted in Section \ref{sec:3.8}, the current experiment pipeline logs the selected feature \emph{count} (feature ratio) for every run but not the identity of \emph{which} features were selected, since only that summary statistic, not the full binary feature mask, is written to the results tables. This was an oversight in the original logging design rather than a deliberate scope decision: the physiological grounding reviewed in Section \ref{sec:2.2}, namely which acoustic measures are expected to reflect parkinsonian dysphonia, would have let the feature subsets CSO converges to be checked for biological plausibility, not just predictive accuracy, which is exactly the kind of evidence a clinician would want before trusting a model's output. Re-running the comparison with the feature mask included in the logged record, which requires no change to any optimizer's search behavior, only to what the experiment harness writes to disk, would let a future revision of this study report which specific acoustic features, and in particular whether nonlinear dysphonia measures such as RPDE, DFA, and PPE are favored over classical jitter/shimmer perturbation measures, are selected most consistently across independent runs.

A second limitation concerns external validity. Both Oxford and Naranjo are public, retrospective, single-recording-protocol datasets with a modest number of subjects (32 and 80, respectively); neither the six-method comparison of Section \ref{sec:4} nor the $\varphi$ sensitivity analysis of Section \ref{sec:sensitivity} has been repeated on a prospectively collected cohort, a different recording device or microphone, a different language, or a clinical population with comorbid conditions that could confound acoustic features. This is a limitation the broader voice-based PD literature shares, not one specific to this study \cite{SedighMalekroodi:2025}, but it means the balanced accuracies in Tables \ref{tab:results-oxford} and \ref{tab:results-naranjo} should be read as evidence that the joint optimization paradigm works on these two datasets, not as an estimate of accuracy on an arbitrary future population.

Third, the $\varphi$ sensitivity analysis of Section \ref{sec:sensitivity} was conducted only for the proposed method's own hyperparameter, not for the five baselines' hyperparameters listed in Table \ref{tab:hyperparameters}, which are instead taken as specified in each baseline's source paper. Those source papers tuned their hyperparameters for their own, typically feature-selection-only, formulations, not for the joint feature-classifier-transfer-function search space introduced in Section \ref{sec:3.3}; it is possible that a baseline re-tuned specifically for this joint search space would narrow some of the gaps reported in Section \ref{sec:5.2}. Reporting the baselines under their published hyperparameters keeps the comparison reproducible and faithful to each source paper, but it is not the same claim as showing every baseline was given an equally thorough hyperparameter search on this exact task.

Fourth, the multi-classifier soft-voting ensemble (Section \ref{sec:3.3}) trains up to five base classifiers per fitness evaluation, and every optimizer in this study spends 3000 such evaluations per run; this is a one-time, offline search cost, not a per-patient inference cost, since only the final selected classifier(s) are needed at deployment time, but it does mean reproducing or extending this comparison, for instance to more datasets, more $\varphi$ values, or more baselines, is markedly more expensive than a feature-selection-only wrapper study would be. Finally, this study addresses a single modality, sustained-phonation acoustic features; PD is a multi-system disorder with established non-vocal digital biomarkers as well \cite{Sun:2024}, and combining voice with gait, tremor, or handwriting features within the same joint optimization paradigm introduced here is a natural direction for future work.
